\section{Conclusion}

We began this work with a mistaken assumption: that SFT fails on hard competitive programming benchmarks because training data lacks correct solutions. APPS contains reference solutions for 85.32\% of competition-level problems. The dataset is not sparse --- the model simply cannot transfer those solutions to held-out hard evaluation problems. That generalization void, not a data void, is where RLEF concentrates its advantage --- though directly confirming the precise mechanism awaits future experimentation (\S\ref{sec:discussion}, L2).

This reframing changes how we should think about why execution feedback helps. RLEF works not because it provides learning signal where SFT has none in the training set, but because it trains on the model's own generated outputs --- learning from what the model can actually attempt rather than from an idealized reference distribution the model cannot yet reach.

\textbf{Summary.} We make three contributions. First, a controlled, reproducible, open-source pipeline for comparing RLEF against SFT across the full benchmark difficulty spectrum using DeepSeek-Coder-7B, APPS, and bigcode-evaluation-harness. Second, empirical confirmation of a statistically significant positive-ordered difficulty-scaling advantage for RLEF over SFT (JT z=+56.10, p$\approx$0), with the largest gains at LiveCodeBench-Hard ($\Delta$=+0.18, proxy from smoke-scale evaluation). Third, a practical null result: fraction-of-tests and binary RLEF rewards produce equivalent performance at APPS training scale (p=0.552), providing a controlled replication of reward formulation convergence on a new model/dataset combination.

\textbf{Future directions.} Immediate next steps: (1) mechanism verification --- re-run h-m2 with max\_new\_tokens$\geq$512; (2) full-scale evaluation --- 4449 samples, 3 epochs, bigcode-harness with saved checkpoints. Medium-term: weighted fraction reward (VeRPO formulation) for multi-test APPS problems; model scale extension to 13B/34B.

The APPS coverage paradox --- correct solutions in the dataset, generalization failure at evaluation --- is a more general phenomenon. As the field scales code LLMs through execution feedback, understanding why models with high-quality training data still fail to generalize to held-out evaluation will be central to making that training effective.
